\section{A functional programming language}
\label{sec:higher-order}

% in this section, we use Haskell
\lstset{
  language=Haskell,
  basicstyle=\ttfamily\small,
  keywordstyle=\color{blue},
  commentstyle=\color{gray}\itshape,
  showstringspaces=false,
  columns=flexible,
}


In this section, we present a functional programming language that defines exactly the polyregular functions. This language grows out of the term based approach from \cref{sec:regular-list-functions}, but thanks to the extra expressive power of polyregular functions, the language can be equipped with programming constructions that make it quite usable.

Before presenting the full programming language, we begin with a minimal adjustment of the term based approach from \cref{sec:regular-list-functions} to cover polyregular functions. This adjustment  is to add one new atomic function, which can be  seen as a kind of for loop that iterates over all splits of an input list.

\begin{definition}[Polyregular terms]\label{def:polyregular-terms}
    A \emph{polyregular term} is a term built using the atomic functions and combinators from \cref{def:regular-terms}, extended one more atomic function:
    \begin{align*}
    A^*  & \xrightarrow{\splitterm} (A^* \times A^*)^* \\
    [a_1,\ldots,a_n] & \mapsto \myunderbrace{[([a_1,\ldots,a_{n}],[]),  ([a_1,\ldots,a_{n-1}],[a_n]), \ldots , ([],[a_1,\ldots,a_{n}])]}{the input list is split into (prefix, suffix) pairs in all possible ways, \\ with the prefix getting shorter as we move to the right}.
\end{align*}
\end{definition}

As we will show in a moment without much difficulty,  the polyregular terms define exactly the polyregular funtions.  We begin, however, with two examples that illustrate the new atomic function.


\begin{myexample} Marked squaring can be defined by the composition
    \begin{align*}
A^* 
\xrightarrow \splitterm 
(A^* \times A^*)^* 
\xrightarrow {(\leftterm^* + \rightterm^*)^*}
((A+A)^* \times (A+A)^*
\xrightarrow {\binconcterm^*} 
(A+A)^{**},
\end{align*}
followed by $\concatterm$ to flatten the output. 
\end{myexample}

\begin{myexample}\label{ex:reverse-from-split}  Reverse  can be implemented as follows.  Suppose that the input list is 
    \begin{align*}
    [a_1,\ldots,a_n] \in A^*.
    \end{align*}
    We first apply $\splitterm$ followed by $\sndterm^*$, yielding a result of the form 
    \begin{align*}
    [[], [a_n], [a_{n-1},a_n], \ldots, [a_1,\ldots,a_n]] \in (A^*)^*.
    \end{align*}
    Next, to each list item  we apply a verson of $\headterm$, which has output $A^* \to A^*$ and returns the singleton of the head if the list is non-empty, and the empty list otherwise. This yields a result of the form
    \begin{align*}
    [[], [a_n], [a_{n-1}], \ldots, [a_1]] \in (A^*)^*.
    \end{align*}
    Applying $\concatterm$ to this result yields the final output.
\end{myexample}

In the exercises we show that also $\blockterm$ becomes redundant in the presence of $\splitterm$, and the group prefix multiplication function can be simplified (although not altogether eliminated).
We now show that the polyregular terms define exactly the polyregular functions. The functions in question are type-to-type functions, with the notion of polyregularity being defined using string representations as in \cref{def:polyregular-functions}. 

\begin{theorem}  \label{thm:polyregular-terms} A type-to-type function is polyregular under string representation  if and only if it can be defined by a polyregular term.
\end{theorem}
\begin{proof}
    The right-to-left implication is easy to show: the combinators are easily seen to preserve polyregularity, the old atomic functions are regular and thus also polyregular, and the new atomic function is easily seen to be polyregular (a for-transducer can have an outermmost loop which marks the comma in the input string which separates the prefix from the suffix). For the left-to-right implication, we use the definition of polyregular functions in terms of compositions of primes from \cref{def:polyregular-functions}; in this case it is enough to show that map squaring can be expressed using $\splitterm$. This is an easy exercise, which we leave to the reader.
\end{proof}





\subsection{Higher-order polyregular terms}
\label{sec:syntax-and-semantics-of-higher-order-terms}

The system from \cref{thm:polyregular-terms} is only a minor variation on the system of combinators from \cref{sec:regular-list-functions}, and is essentially a warm-up for the real contribution of this section, which is a kind of functional programming language. 

Functional programming languages are based on the $\lambda$-calculus, and the basic computational mechanism in the $\lambda$-calculus is substitution, which replaces a variable by its argument. Since a variable can appear multiple times, such substitution can lead to quadratic growth. For example, the squaring operation from \cref{exer:squaring-continuous} can be implemented by the following  code, which replaces each list element by the list itself:
\begin{align*}
\lambda x. \concatterm\ ( \mapterm\ (\lambda y. x)\  x)
\end{align*}
Because of the non-linear outputs, we could not have substitution in the system of combinators from \cref{sec:regular-list-functions}, which was designed to capture the regular functions that have linear growth. However, in the polyregular setting, we can embrace substitution. This is what we do in this section: we define a functional programming language which defines exactly the polyregular functions. 

Like essentially all functional programming languages, our language will have higher-order functions, so that  functions can input other functions, or even more complicated objects. For example, a legitimate type in our language will be 
\begin{align*}
(\onetype^* \to (\onetype + \onetype)) \to \onetype + \onetype,
\end{align*}
which describes a function that inputs a subset of the natural numbers and outputs a Boolean. The functions can be themeselves nested with the other type constructors, e.g.~we have a type for lists of functions as in the above type.  The corresonding set of types is defined below.





\begin{definition}
    [Higher-order types] \label{def:higher-order-types}  The \emph{higher-order types} are defined by extending the grammar of types in \cref{def:types}  an arrow constructor $A \to B$. 
\end{definition}

As was the case in \cref{sec:regular-list-functions}, each type represents a set. For the arrow constructor $A \to B$, the represented set consists of total functions from the set represented by $A$ to the set represented by $B$. By abuse of notation, we use the same symbol for a type and the set that it represents.


We now describe the syntax of our programming language. Essentially, this is the extension of the polyregular terms with the mechanics of the  $\lambda$-calculus. We use a variant of Church style typing, where each variable comes hard-wired with a type, and therefore each term has a unique type, which can be inferred by a bottom-up analysis. For the rest of this section, we assume that there is some set of variables, each one with an associated higher-order type, such that each higher-order type has infinitely many variables. 
\begin{definition}
    [Higher-order polyregular terms] 
    \label{def:higher-order-terms}
    The \emph{higher-order  polyregular terms} and their associated types  are defined by induction as follows (for brevity, in the definition we use the name ``term'' meaning higher-order polyregular term):
        \begin{description}
            \item[variable:]        each variable is a term of corresponding type;
            \item [application] if $f : A \to B$ and $a : A$ are terms, then $f(a) : B$ is a term;
        \item[atomic functions:] we keep the    atomic polyregular functions\footnote{In fact, only the subset listed in \cref{lem:fewer-atomic} is enough.}, which are  now also allowed for higher-order types, 
            % \begin{align*}
            % \fstterm : &\  A \times B  \to A \\
            % \sndterm : &\  A \times B  \to B \\
            % \unconsterm : &\  A^* \to \onetype + A \times A^*\\
            % \concatterm : &\ A^* \to A \\
            % \splitterm : &\ A^* \to (A^* \times A^*)^*  \\
            % \multterm : &\  G^* \to G 
            % \end{align*}
        we add two extra atomic functions:
        \begin{align*}
        \mapterm & : (A \to B) \to A^* \to B^* \\
        \mathtt{either}  & : (A \to C) \to (B \to C) \to (A+B) \to C
        \end{align*}
                \item[unit:] there is a unit term $1 : \onetype$;
            \item[pairs:] if $a : A$ and $b : B$ are terms, then so is $(a,b) : A \times B$;
            % \item [co-pairs:]
            % \ 
            % \begin{enumerate}
            %     \item if $a  : A$ is a term and $B$ is a type, then  $\leftterm a : A + B$ is term;
            %     \item if $b  : B$ is a term and $A$ is a type, then  $\rightterm b : A + B$ is term;
            % \end{enumerate}
            \item [lists] if $a_1,\ldots,a_n : A$ are terms, then $[a_1,\ldots,a_n] : A^*$ is a term;
            \item [abstraction] if $b : B$ is a term and $x : A$ is a variable, then $\lambda x. b : A \to B$ is a term.
        \end{description}
\end{definition}

 Observe that we no longer need the four combinators from \cref{sec:regular-list-functions}, since composition is a special case of abstraction, map and co-pairing have become  atomic functions, and pairing is now a term constructor. We think of the last four items in the above definition, i.e.~unit, pairing, lists and abstraction, as term constructors for the corresponding type constructors. We also think of the atomic functions $\leftterm$ and $\rightterm$ as term constructors for the co-product type, even though they are not explicitly listed since they happen to be already present as atomic polyregular functions. 
  
 \paragraph*{Semantics.} If a higher-order polyregular  term of type $A$ is closed, i.e.~has no free variables, then its semantics is an element of  (the set represented by) $A$. To define the semantics formally, we proceed by induction on the structure of the term, which involves taking care of terms with free variables. 
  For a term of type $A$ which has  has free variables of types $A_1,\ldots,A_n$,  its semantics is the  function  of type 
\begin{align*}
 {A_1} \times \cdots \times  {A_n} \to A
\end{align*}
which inputs values for the free variables and returns the result of the term. This function is defined in the expected way by induction on the structure of the term.  


\paragraph*{Redundant atomic functions.} Some of the other atomic functions  become redundant in the presence of the new features of higher-order polyregular terms, such as: 
         \begin{align*}
        \idterm  &=  \lambda x.x \\
        \diagterm  & = \lambda x. (x,x) \\
        \codiagterm & = \either (\lambda x.x) (\lambda x.x) \\
        \distrterm & = \lambda x.\  \either 
        (\lambda y. \leftterm\ (\fstterm\ x, y))\  
        (\lambda y. \rightterm\ (\fstterm\ x, y))\ 
        (\sndterm\ x) \\
        \consterm  & = \either (\lambda y. [])\ 
        (\lambda y. \concatterm [ [\fstterm y], \sndterm y]] ) 
    \end{align*}
We have also removed the reverse function in Example~\ref{ex:reverse-from-split}, and the block function in~\cref{exer:polyregular-terms-eliminate-block}. Finally, in 
\cref{exer:polyregular-terms-simplify-group}, we show that the group prefix multiplication can be replaced by a simpler function 
\begin{align*}
G^* \xrightarrow {\multterm} G
\end{align*}
which computes the group product of all elements in the input list (in other words, this is the last letter of the group prefix multiplication, or the group identity if the input list is empty).
The following lemma summarises these results, by listing the remaining atomic functions. 


\begin{lemma}\label{lem:fewer-atomic}
    The expressive power the higher-order polyregular terms is  not changed if the set of atomic functions is restricted to the   co-projections (which are now called term constructors) and the following atomic functions:
    \begin{align*}
            \fstterm : &\  A \times B  \to A \\
            \sndterm : &\  A \times B  \to B \\
            \unconsterm : &\  A^* \to \onetype + A \times A^*\\
            \concatterm : &\ A^* \to A \\
            \splitterm : &\ A^* \to (A^* \times A^*)^*  \\
            \multterm : &\  G^* \to G  \\
        \mapterm  : &\   (A \to B) \to A^* \to B^* \\
        \mathtt{either}  : &\  (A \to C) \to (B \to C) \to (A+B) \to C
        \end{align*}
\end{lemma}



\subsection{Expressive completeness}
\label{sec:expressive-completeness}
We now state the main result of this section, which is that the higher-order polyregular  terms define exactly the polyregular functions. To compare the two notions, we need to talk about string-to-string functions, or the slightly more general case of functions between types that are arrow-free, since only such functions have a defined notion of  polyregularity.





\begin{theorem}\label{thm:higher-order-terms}
    Consider a function 
    \begin{align*}
    A \xrightarrow f B
    \end{align*}
    where the types $A$ and $B$ are arrow-free, i.e.~constructed without using the arrow type constructor. Then $f$ is polyregular if and only if it can be defined by a higher-order term.
\end{theorem}

    The left-to-right implication is immediate, since the higher-order polyregular terms were designed as an extension of the polyregular terms from \cref{def:polyregular-terms}. The hard direction is the converse implication, which says that the higher-order features do not add expressive power, at least as long as we define functions between arrow-free types. 
    

    For the rest of this section, we will say ``term'' instead of ``higher-order polyregular term'' for brevity. We also assume that the set of atomic functions is the reduced set from \cref{lem:fewer-atomic}. 

        \newcommand{\reduces}{\rightsquigarrow}

            \begin{figure}
                \begin{align*}
            \fstterm (a,b) & \reduces a\\
            \sndterm (a,b) & \reduces b\\
            \unconsterm [] & \reduces \leftterm 1\\
            \unconsterm [a_1,\ldots,a_n] & \reduces \rightterm (a_1, [a_2,\ldots,a_n])\\
            \concatterm [\ell_1,\ldots,\ell_n] & \reduces \ell \quad \text{(see the caption for restrictions)}  \\
            \splitterm [a_1,\ldots,a_n] & \reduces [([a_1,\ldots,a_n],[]),  ([a_1,\ldots,a_{n-1}],[a_n]), \ldots , ([],[a_1,\ldots,a_{n}])] \\
            \multterm [g_1,\ldots,g_n] & \reduces g
            \quad \text{(see the caption for restrictions)}  \\
        \mapterm f [a_1,\ldots,a_n] & \reduces [f(a_1),\ldots,f(a_n)]\\
        \either f g \ \leftterm\ a & \reduces f(a)\\
        \either f g \ \rightterm\ b & \reduces g(b)
        \end{align*}
    \caption{
    \label{fig:reduction-rules}    
    Reduction rules for the atomic functions. Let us explain the restrictions for the rules for $\concatterm$ and $\multterm$.
        In the reduction rule for $\concatterm$, the terms $\ell_1,\ldots,\ell_n$ must be list terms, which means that each one is a list of terms (but the list elements are not restricted in any way). The ouput  term $\ell$ is then their concatenation.  In the reduction rule for $\multterm$, the terms $g_1,\ldots,g_n$ must be group values, which means that they  must be constructed by the term constructors for $\onetype$, products and co-products (in particular, each group element has a unique representation as a term). The output term $g$ is then required to be the value the represents the group multiplication of the inputs.
    }
    \end{figure}


    \paragraph*{Reduction.} Our proof of the difficult direction will be based on evaluating terms via $\beta$-reduction. The main reduction rule is 
    \begin{align*}
    (\lambda x. a)(b) \reduces  \myunderbrace{a[x \mapsto b]}{term $a$ with free occurrences of variable $x$ replaced by term $b$},
    \end{align*}
    which substitutes the variable $x$ for the argument. Apart from this rule, we have reductio rules for the atomic functions, which are listed in \cref{fig:reduction-rules}. Finally, all reduction rules can be applied to subterms. This gives us a notion of one-step reduction between terms, which we denote by $a \reduces b$. It is easy to see that one-step reduction changes neither the type nor the semantics of a term. A term is said to be in \emph{normal form} if no reduction rules can be applied to it. A classical result of the $\lambda$-calculus is that the normal form is unique if it exists, but we will not need to prove this. 

    \paragraph*{String representation.} To implement reduction by polyregular functions, we will use string representations of terms. We will only be interest in terms which are closed. In order to avoid special letters for variables, we use the following convention, which is based on De Bruijn indices: in an abstraction we do not give the name to a variable but only its type, and when referencing a variable we only say how many steps up in the syntax tree the variable is defined. The number of steps is represented by repetitions of an arrow $\uparrow$. Here are two examples of this representation:
    \begin{align*}
   \lambda x : A. x \quad & \text{is represented by } \lambda : A.  \uparrow\\
   \lambda x: A.\ \lambda y: B.\ x \quad & \text{is represented by } \lambda : A. \lambda : B.\  \uparrow \uparrow 
    \end{align*}
    Under this representation, every closed term can be written as an alphabet which contains the eight lettters used in string representations of regular terms from \cref{sec:regular-list-functions}, plus the following letter:
    \begin{align*}
    \myunderbrace{\lambda\ :}{a single \\ letter}
    \quad \uparrow \quad  \myunderbrace{\onetype \quad  + \quad \times \quad * \quad \rightarrow }{needed to write types}.
    \end{align*}
    We will call this alphabet $F$, because it is finite and has fifteen letters (which is $F$ in hexadecimal). Summing up, each closed term can be represented by a string over the alphabet $F$, and we will call this string the \emph{string representation} of the term. This representation is unique, as long as we write the types in a unique way, which can be done by always using parentheses. 
    
    \paragraph*{Parsing.} In order to parse closed terms using a finite automaton, we will impose  bounds on the types that appear in a term (i.e.~are used by subterms of it) and the height of the term (which is the height of the syntax tree, with the height of a list term being one plus the maximal height of the list elements).  Here is a formal definition.

    \begin{definition}[Resource bound]\label{def:resource-bound}
        A \emph{resource bound} is a pair consisting of a finite set $\Aa$ of types and a height bound $h \in \set{1,2,\ldots}$. A closed term is said to be \emph{consistent} with the resource bound  if its height is at most $h$ and all types that appear in the term are in $\Aa$.
    \end{definition}




    \begin{lemma}\label{claim:parsing-is-regular-with-resource-bounds}
        For every resource bound $(\Aa,h)$, the following is a regular language
        \begin{align*}
        \setbuild{ w \in F^*}{$w$ is the string representation of a term \\ that  is consistent with the resource bound $(\Aa,h)$}.
        \end{align*}
    \end{lemma}
    \begin{proof}
        Straightforward induction on the height.
    \end{proof}

    So far, the resource bounds are just a syntactic device which is used to enable parsing. However, they will also have an important semantic role in our approach, since resource bounds will allow us to control the reduction of terms.

    \paragraph*{Reduction strategy.}  A \emph{redex} in a term is defined to be a node in the syntax tree where a reduction can be applied.
    A \emph{reduction strategy} is a function which assigns to each term a set of non-overlapping redexes in it. A reduction strategy is can be used to transform terms: we take all redexes given by the reduction strategy, and reduce them in parallel.  

    \begin{lemma}\label{lem:reduction-strategy-lemma}
    There is a reduction strategy which has the following properties:
    \begin{enumerate}
        \item \textbf{Termination.} For every resource bound $(\Aa,h)$ there is some  $k \in \set{1,2,\ldots}$ such every  term consistent with $(\Aa,h)$ normalises after at most $k$ steps of the strategy.
        \item \textbf{One-step.}  For every resource bound $(\Aa,h)$, there is some polyregular function $f$ such that the following diagram commutes:
\[
\begin{tikzcd}[column sep=3cm]
\text{closed terms consistent with $(\Aa,h)$}
\ar[d, "\text{string representation}"']
\ar[r, "\text{one step of the strategy}"] 
&
\text{closed terms} 
\ar[d, "\text{string representation}"]
\\
F^*
\ar[r,"f"'] & F^*.
\end{tikzcd}
\]
    \end{enumerate}
    \end{lemma}


        Before proving the lemma, we use it complete the proof of \cref{thm:higher-order-terms}.
        
\begin{proof}[Proof of \cref{thm:higher-order-terms} using \cref{lem:reduction-strategy-lemma}]
    We have already shown that every polyregular function can be defined by a higher-order term. Let us now prove the converse implication.          The essential idea is to compute the output $f(a)$ by first writing a string representation of the term $f(a)$, and then reducing this term until normal form is reached. The reduction strategy from \cref{lem:reduction-strategy-lemma}  is used to prove the following corollary, which says that terms of given resource bounds can be normalised using a polyregular function.
    \begin{corollary}\label{cor:normalisation-of-terms-with-resource-bounds}
        For every resource bound there is a polyregular function, which computes the normal form for terms consistent with the resource bound (both input and output are given in string representations).
    \end{corollary}
    \begin{proof}
        The termination item in \cref{lem:reduction-strategy-lemma} gives a bound on the number of reduction steps needed to reach normal form. The one-step item in the lemma says that each step can be computed by a polyregular function. There is one minor item that needs to be addressed: the resource bounds might change during the reduction. However, each reduction step can at most double the height of a term (up to an additive constant), and the only new types that can appear are subtypes of types that were already in the term. Therefore, if $k$ is the bound from the termination lemma, then there exist resource bounds 
        \begin{align*}
        (\Aa_1,h_1),\ldots, (\Aa_k,h_k),
        \end{align*}
        such before applying the $i$-th step of the  reduction strategy, the term is consistent with $(\Aa_i,h_i)$, and after $k$ steps it is in normal form. Therefore, each step can be computed by a polyregular function, and the composition of these $k$ functions is also polyregular.
    \end{proof}

Using the above corollary, we can now complete the proof of \cref{thm:higher-order-terms}. Indeed, suppose that $f$ is a term of type $A \to B$ where both $A$ and $B$ are arrow-free. Because the input and output types are arrow-free, we will be able to use the following observation about normal forms of terms.

\begin{lemma}\label{lem:normal-forms-of-arrow-free-terms}
    Consider a closed term whose type is arrow-free. Then this term is in normal form if and only if it is built using the value constructors $(a,b)$, $[a_1,\ldots,a_n]$, $\leftterm a$ and $\rightterm b$, and $1$.
\end{lemma}
\begin{proof}
    An inspection of the shapes of closed terms in normal form. 
\end{proof}

Suppose that we have an input $a \in A$. We can view the application $f(a)$ as a term, in which the source code of $f$ is followed by a term that describes $a$. This application has resource bounds that depend only on the term $f$ and the input type $A$. Therefore, we can apply \cref{cor:normalisation-of-terms-with-resource-bounds} to compute the normal form of the term $f(a)$ using a polyregular function.  By \cref{lem:normal-forms-of-arrow-free-terms}, this normal form is the same a string representation of the output value. Hence, the function $f$ is polyregular under string representations, which completes the proof of \cref{thm:higher-order-terms}. 
\end{proof}
It remains to prove \cref{lem:reduction-strategy-lemma}.


\newcommand{\ordof}{\mathsf{ord}}
\newcommand{\steps}{\mathsf{steps}}
\newcommand{\Maj}{\mathcal M}

\begin{proof}[Proof of \cref{lem:reduction-strategy-lemma}]
    We use the \emph{outermost} strategy: the redexes selected in a term are those which are not proper subterms of other redexes. Of two nested redexes the inner one is not outermost, and therefore the selected redexes are pairwise non-overlapping, as required of a reduction strategy. We write
    \begin{align*}
    t \Rightarrow t'
    \end{align*}
    for one step of this strategy, i.e.~for the result of contracting all outermost redexes of $t$ in parallel. Since a term which is not in normal form contains a redex, and hence also an outermost one, the step $t \Rightarrow t'$ is available precisely when $t$ is not in normal form. The strategy is deterministic, and therefore each term $t$ has a unique maximal sequence
    \begin{align*}
    t = t_0 \Rightarrow t_1 \Rightarrow t_2 \Rightarrow \cdots;
    \end{align*}
    we write $\steps(t) \in \set{0,1,\ldots} \cup \set \infty$ for its length. In this notation, the termination item in the lemma says that $\steps(t) \le k$ holds for all terms $t$ that are consistent with the resource bound. We prove the two items of the lemma separately.

    \paragraph*{Termination.} Let $\Bb$ be the closure of $\Aa$ under subtypes; this is again a finite set of types. As we observed in the proof of \cref{cor:normalisation-of-terms-with-resource-bounds}, a reduction step introduces no types other than subtypes of types that are already present. Therefore, if $t$ is consistent with the resource bound $(\Aa,h)$, then all types appearing in $t$, and in all of its reducts, belong to $\Bb$. The height, on the other hand, will grow during the reduction, and the lengths of the lists inside $t$ are not bounded by the resource bound at all. This is why the argument below is arranged so that neither of these two quantities enters the final estimate.

    The proof is the majorant method, i.e.~a step-counting version of Tait's reducibility argument. To each term we assign a \emph{majorant}, which bounds the number of steps needed to normalise the term, and also, hereditarily, the same information about everything that can be extracted from its normal form. For a type $A$, the set $\Maj(A)$ of majorants of type $A$ is defined by induction on $A$, together with a \emph{cost} function $\Maj(A) \to \Nat$ which is written $\mu \mapsto |\mu|$:
    \begin{align*}
    \Maj(\onetype) &= \Nat  &  |m| &= m\\
    \Maj(A \times B) = \Maj(A + B) &= \Maj(A) \times \Maj(B) &  |(\mu,\nu)| &= \max(|\mu|,|\nu|)\\
    \Maj(A^*) &= \Maj(A)  & & \\
    \Maj(A \to B) &= \Nat \times [\Maj(A) \to \Maj(B)] & |(m,F)| &= m,
    \end{align*}
    where $[\Maj(A) \to \Maj(B)]$ denotes the set of monotone functions. Each $\Maj(A)$ is ordered coordinatewise, and pointwise in the case of the function space; a straightforward induction on the type shows that every two majorants $\mu,\nu \in \Maj(A)$ have a least upper bound $\mu \vee \nu$, and that the cost function is monotone. We also write $\mu + r$ for the majorant obtained from $\mu$ by adding $r$ to each of the numbers that appear in it; this operation is monotone and satisfies $|\mu + r| = |\mu| + r$.

    We now define what it means for a majorant $\mu \in \Maj(A)$ to \emph{majorise} a term $t : A$, which need not be closed. The definition is by induction on the type. In every case we require
    \begin{align*}
    \steps(t) \le |\mu|,
    \end{align*}
    which in particular asserts that $t$ has a normal form. Apart from this, we require the following for \emph{every} term $s$ that is reachable from $t$ by steps of the strategy, i.e.~such that $t \Rightarrow^* s$:
    \begin{description}
        \item[unit] if $A = \onetype$, nothing more;
        \item[product] if $A = B \times C$ and $\mu = (\mu_B,\mu_C)$ and $s = (b,c)$, then $b$ is majorised by $\mu_B$ and $c$ by $\mu_C$;
        \item[co-product] if $A = B + C$ and $\mu = (\mu_B,\mu_C)$ and $s = \leftterm\ b$, then $b$ is majorised by $\mu_B$, and symmetrically for $\rightterm$;
        \item[list] if $A = B^*$ and $s = [b_1,\ldots,b_p]$, then \emph{every} entry $b_i$ is majorised by $\mu$;
        \item[arrow] if $A = B \to C$ and $\mu = (m,F)$, then for every term $u : B$ and every majorant $\nu$ of $u$, the term $s\ u$ is majorised by $F(\nu)$.
    \end{description}
    Quantifying over the reducts $s$, rather than only over $t$ itself, costs us nothing and buys the closure property that we need below.
    The conditions on $v$ are stated conditionally, so that they are vacuous for the terms whose normal form is headed by a free variable; such terms are called \emph{neutral} below. The list item is the heart of the matter: a list is charged the maximum of the majorants of its entries, and not their sum. This is what the parallel strategy buys us, and it is the reason why the final bound will not depend on the lengths of the lists.

    We record three properties of this notion.
    \begin{enumerate}
        \item \label{item:majorant-monotone} \textbf{Monotonicity.} If $\mu$ majorises $t$ and $\mu \le \mu'$, then $\mu'$ majorises $t$. This is a routine induction on the type, using monotonicity of the cost function and of the majorants of arrow type.
        \item \label{item:majorant-reduction} \textbf{Forward closure.} If $t \Rightarrow t'$ and $\mu$ majorises $t$, then $\mu$ majorises $t'$. This holds by construction: $\steps(t') = \steps(t) - 1 \le |\mu|$ because the strategy is deterministic, and every reduct of $t'$ is a reduct of $t$, so the remaining conditions are inherited verbatim.
        \item \label{item:majorant-neutral} \textbf{Neutral terms.} If $x$ is a variable, and the terms $u_1,\ldots,u_r$ are majorised by $\nu_1,\ldots,\nu_r$, then $x\ u_1 \cdots u_r$ is majorised by any majorant whose cost is at least $|\nu_1| + \cdots + |\nu_r|$ and which is above the majorants required by the arrow item; such a majorant exists, and can be chosen monotonically in the $\nu_i$, since a term whose head is a variable stays neutral under reduction, so that the conditions in the product, co-product and list items are vacuous for it. In particular, every variable of type $A$ is majorised by some $\iota_A \in \Maj(A)$.
    \end{enumerate}
    What we do \emph{not} get for free is the converse of item~\ref{item:majorant-reduction}, the head expansion property: if $\mu$ majorises $b[x \mapsto u]$, then $\mu+1$ majorises $(\lambda x.b)\, u$. \issue{This is the one step of the termination argument that is not routine, and it should be checked carefully before the proof goes into the book. The difficulty is the interaction between the parallel strategy and substitution: the reducts of $(\lambda x.b)u$ are the reducts of $b[x \mapsto u]$, which is what we want, but in the arrow item we must compare $((\lambda x.b)u)\,w$ with $(b[x \mapsto u])\,w$, and one step of the strategy reduces the argument $w$ at the same time as the head redex. The standard remedy is to prove the substitution lemma for parallel reduction ($b \Rightarrow b'$ and $u \Rightarrow u'$ imply $b[x \mapsto u] \Rightarrow b'[x \mapsto u']$) and to sequentialise application spines in the strategy, so that in an application whose function part is not yet normal only the function part is reduced. Everything else in this proof is independent of how this point is settled.} We take this property for granted below; it is the standard bookkeeping of the majorant method, and it is the only place where the parallel character of the strategy has to be reconciled with substitution.

    \begin{claim}\label{claim:atomic-functions-majorised}
        Every atomic function is majorised by some majorant.
    \end{claim}
    \begin{proof}
        We give two representative cases; the remaining atomic functions are treated in the same way. Both cases use the following observation. Suppose that a term of the form $g\ u$ is not a redex, and will not become one until the reduct of $u$ acquires the shape required by the reduction rule for $g$. As long as this has not happened, the outermost redexes of $g\ u$ are exactly the outermost redexes of $u$, and therefore the first steps of the strategy applied to $g\ u$ are the steps of the strategy applied to $u$. In particular, if $u$ is majorised by $\mu$ and the required shape appears after $j \le |\mu|$ steps, then, by item~\ref{item:majorant-reduction} above, the parts of the $j$-th reduct of $u$ are still majorised by the corresponding parts of $\mu$, and they normalise in at most $|\mu| - j$ further steps.

        Consider first $\fstterm : A \times B \to A$, and let $u$ be majorised by $\mu = (\mu_A,\mu_B)$. If the normal form of $u$ is not a pair, then $\fstterm\ u$ is neutral, and item~\ref{item:majorant-neutral} applies. Otherwise, after $j \le |\mu|$ steps the reduct of $u$ becomes a pair $(a,b)$, where $a$ is majorised by $\mu_A$ and satisfies $\steps(a) \le |\mu| - j$. The next step contracts $\fstterm\ (a,b)$ into $a$, and the reduction continues inside $a$. Altogether
        \begin{align*}
        \steps(\fstterm\ u) \ \le \ j + 1 + (|\mu| - j) \ = \ |\mu| + 1,
        \end{align*}
        and $\fstterm\ u$ is majorised by $\mu_A + (|\mu|+1)$, joined with the majorant supplied by item~\ref{item:majorant-neutral}. Since this assignment is monotone in $\mu$, it is the function part of a majorant of $\fstterm$.

        Consider now $\mapterm : (A \to B) \to A^* \to B^*$, let $f$ be majorised by $(m,F)$ and let $\ell$ be majorised by $\rho$. Reasoning as above, after $j \le |\rho|$ steps the reduct of $\ell$ becomes a list $[a_1,\ldots,a_q]$ whose entries are all majorised by $\rho$, while $f$ has become a term $f'$ that is still majorised by $(m,F)$. One further step contracts the redex into
        \begin{align*}
        [f'(a_1),\ldots,f'(a_q)],
        \end{align*}
        and by the arrow item every entry of this list is majorised by $F(\rho)$. The entries are now reduced in parallel, so the whole list is majorised by $F(\rho) + (|\rho|+1)$. Note that the number $q$ of entries has played no role whatsoever: this is the list item at work. Since the assignment $((m,F),\rho) \mapsto F(\rho) + (|\rho|+1)$ is monotone in both arguments, it gives the required majorant of $\mapterm$.
    \end{proof}

    \begin{claim}\label{claim:every-term-majorised}
        For every term $t : A$ with free variables $x_1 : A_1, \ldots, x_q : A_q$ there is a monotone function
        \begin{align*}
        \Phi_t : \Maj(A_1) \times \cdots \times \Maj(A_q) \to \Maj(A)
        \end{align*}
        such that $\Phi_t(\nu_1,\ldots,\nu_q)$ majorises $t[x_1 \mapsto u_1, \ldots, x_q \mapsto u_q]$ whenever each $u_i$ is majorised by $\nu_i$.
    \end{claim}
    \begin{proof}
        Induction on the structure of $t$. For a variable, $\Phi_t$ is the corresponding projection. For an atomic function, $\Phi_t$ is the constant given by \cref{claim:atomic-functions-majorised}. For a pair $(a,b)$, we take the pair of the majorants given by the induction assumption, which is legitimate because the outermost redexes of $(a,b)$ are those of $a$ together with those of $b$, so that the two components are normalised in parallel. For a list $[a_1,\ldots,a_p]$, we take the join of the majorants of the entries, for the same reason. For an application $f(a)$, we apply the function part of the majorant of $f$ to the majorant of $a$, which is exactly what the arrow item provides. Finally, consider an abstraction $\lambda x.b$, where $x : B$ and $b : C$, and let $\Psi$ be the monotone function given for $b$ by the induction assumption, with $x$ among its arguments. For a majorant $\nu$ of a term $u : B$ we have $(\lambda x. b) u \Rightarrow b[x \mapsto u]$, so the head expansion property stated above shows that $(\lambda x.b)u$ is majorised by $\Psi(\ldots,\nu)+1$. We may therefore take
        \begin{align*}
        \Phi_{\lambda x. b}(\ldots) \ = \ \big(\ |\Psi(\ldots,\iota_B)| \ , \ \nu \mapsto \Psi(\ldots,\nu)+1\ \big),
        \end{align*}
        where the first coordinate is justified by the fact that the reduction of $\lambda x.b$ is the reduction of its body $b$ with $x$ left as a free variable, and a variable is majorised by $\iota_B$.
    \end{proof}

    It remains to observe that, for terms of bounded height with types in $\Bb$, the majorants given by \cref{claim:every-term-majorised} can be chosen uniformly. For $j \in \set{0,1,\ldots}$ and types $A_1,\ldots,A_q,A$ in $\Bb$, define
    \begin{align*}
    \Phi^j_{A_1,\ldots,A_q,A} : \Maj(A_1) \times \cdots \times \Maj(A_q) \to \Maj(A)
    \end{align*}
    by induction on $j$, as the join of the finitely many monotone functions that are obtained by applying to the functions $\Phi^{j-1}$ one of the constructions used in the proof of \cref{claim:every-term-majorised}, together with $\Phi^{j-1}$ itself. Here it is important that the number of these constructions is finite, which is the case because $\Bb$ is finite, and because for lists the construction is a join, so that all list lengths $p$ give the same function. Since a finite join of monotone functions is monotone, and $\Phi^{j-1} \le \Phi^{j}$, an induction on $j$ shows that
    \begin{align*}
    \Phi_t \ \le \ \Phi^j_{A_1,\ldots,A_q,A}
    \end{align*}
    holds for every term $t$ of height at most $j$ whose types are in $\Bb$. Apply this to a closed term $t$ that is consistent with the resource bound $(\Aa,h)$: such a term has height at most $h$ and types in $\Bb$, and therefore it is majorised by $\Phi^h_A \in \Maj(A)$, where $A \in \Bb$ is its type. Consequently
    \begin{align*}
    \steps(t) \ \le \ k \qquad \text{where} \qquad k \ = \ \max \setbuild{|\Phi^h_A|}{$A \in \Bb$},
    \end{align*}
    and the maximum is over a finite set. This completes the proof of the termination item. Let us stress where the two unbounded parameters went: the lengths of the lists disappeared in the list item, and the growth of the height during the reduction was never used, because the majorant of a term is computed once, from its original syntax, and is preserved by reduction.

    \paragraph*{One step.} Fix a resource bound $(\Aa,h)$. We describe a polyregular function which, given the string representation of a closed term $t$ that is consistent with $(\Aa,h)$, outputs the string representation of the term $t'$ with $t \Rightarrow t'$. On inputs which are not such representations the function may do anything, since the diagram in the lemma constrains it only on the representations; by \cref{claim:parsing-is-regular-with-resource-bounds} these form a regular language, so we may as well output the empty string on the remaining inputs. We use for-transducers, which compute exactly the polyregular functions by \cref{thm:for-transducers-are-polyregular}.

    The first observation is that a bounded amount of information about the syntax tree is available in a regular way. Since the height of $t$ is at most $h$, the list of ancestors of a position in the syntax tree has length at most $h$, and therefore it is a bounded object which can be stored in the state of a finite automaton. Combining this with \cref{claim:parsing-is-regular-with-resource-bounds}, we see that there is a rational function which annotates each letter of the input with: the list of the ancestors of the corresponding node together with their types; the number of $\lambda$'s among these ancestors; and a flag which says whether the position begins an outermost redex and, if so, which reduction rule applies to it. The only point that deserves a comment is that being a redex is a regular property of a node: for all rules except those for $\concatterm$ and $\multterm$ it is determined by the first letter of the argument, and for these two rules it amounts to checking that the argument is a list of list terms, respectively a list of group values, which is again a regular property. Being \emph{outermost} then means that no ancestor is a redex, and the ancestors are bounded in number. Since rational functions are polyregular, and polyregular functions are closed under composition, we may assume from now on that the input comes with this annotation.

    The second observation is that the selected redexes occupy pairwise disjoint intervals of the input, since they are pairwise non-overlapping subtrees. Consequently, the output is obtained from the input by replacing, for each selected redex, the interval that it occupies by the string representation of its contractum, and leaving all remaining letters unchanged. This is implemented by a for-transducer whose outer loop ranges over the positions $x$ of the input. If $x$ lies inside a selected redex but does not begin one, nothing is output. If $x$ lies outside all selected redexes, the letter at $x$ is output. If $x$ begins a selected redex, we output the contractum of that redex, using the annotation to see which rule applies. It remains to describe, for each rule, how the contractum is produced; in each case this is done by loops nested inside the outer loop, so the result is again a for-transducer.

    For the rules of $\fstterm$, $\sndterm$, $\unconsterm$ and $\either$, the contractum is obtained from the redex by copying one or two of its parts and inserting a bounded number of letters, which is done by a single inner loop that ranges over the interval of the redex and outputs, or does not output, the current letter, depending on the annotation. The rule for $\concatterm$ is of the same kind: the inner loop copies the interval of the argument, suppressing the brackets of the inner lists. The rule for $\multterm$ is slightly different, because the output does not appear in the input: here an inner loop ranges over the entries of the argument list and, using Boolean variables to store an element of the finite group $G$, computes the product of the entries; a bounded case distinction on the resulting element then outputs its string representation, which is a fixed string. The rule for $\mapterm$ needs two nested loops: the outer of the two ranges over the entries $a_i$ of the argument list, and for each of them the inner one first copies the interval of $f$ and then the interval of $a_i$. The rule for $\splitterm$ is handled in the same way, and is the reason why a single loop would not suffice: the outer of the two loops ranges over the positions at which the input list is cut, and the inner one copies first the prefix and then the suffix determined by the cut. This is exactly the pattern of the for-transducer for marked squaring in Example~\ref{ex:marked-squaring-for-transducer}, and it is where the quadratic growth of polyregular functions is used.

    We are left with the $\beta$-rule, which is the only rule that involves substitution. Let the selected redex be $(\lambda : A.\ a)\ b$. Its contractum is the term $a$ in which every occurrence of the variable bound by this $\lambda$ is replaced by $b$. Because we use De Bruijn indices, this requires some bounded arithmetic on the indices, which are all at most $h$ and hence can be kept in the state:
    \begin{itemize}
        \item an occurrence $\uparrow^r$ inside $a$ refers to the variable bound by our $\lambda$ exactly when $r$ is one plus the number of $\lambda$'s strictly between the occurrence and our $\lambda$, a number which the annotation provides;
        \item an occurrence with a larger $r$ refers to a $\lambda$ above the redex, and since our $\lambda$ disappears, such an occurrence has to be rewritten as $\uparrow^{r-1}$;
        \item when a copy of $b$ is placed under some number $s \le h$ of additional $\lambda$'s, the indices of the free variables of $b$ have to be increased by $s$.
    \end{itemize}
    All three operations replace a letter by a letter, with the replacement determined by the annotation, so they are performed on the fly. The contractum is therefore produced by two nested loops: the outer one ranges over the interval of $a$ and outputs the current letter, unless this letter begins an occurrence of the bound variable, in which case the inner one copies the interval of $b$ with the indices shifted as described above.

    Summing up, one step of the strategy is computed by a for-transducer, whose loops are nested to depth three, composed with a rational annotating function. By \cref{thm:for-transducers-are-polyregular} this function is polyregular, which completes the proof of the one-step item, and with it the proof of the lemma.
\end{proof}
